\scanpage{61}
\begin{tabular}{c@{\qquad}l}
U N I V E R S I T E & MONTPELLIER, le 31 MAI 1979\\
DES SCIENCES ET TECHNIQUES &\\
DU LANGUEDOC &\\
\texttt{------------} &
\end{tabular}

Service de la Scolarité\\
MR / GJ\\
Diffusion : Liste A\\
N° : 3159

\uline{AVIS DE SOUTENANCE D'UNE THESE DE DOCTEUR INGENIEUR}

\uline{Chimie industrielle - l'Eau et les Nuisances}

\uline{Nom et prénom du candidat} : B L A N C H E T O N\quad J e a n - P a u l

\uline{Titre de la thèse} : Analyse du fonctionnement d'un étang-solaire par modélisation à partir de données recueillies sur un étang expérimental à Palavas.

\uline{Date, heure et lieu de soutenance prévus} : 12 JUIN 1979 à 14 Heures 30 mn\\
Salle de cours n° 9 - U.S.T.L.

\uline{Nom du Directeur de thèse} : Monsieur le Professeur BOCQUILLON Claude.

\uline{Désignation du laboratoire où la thèse est préparée:}\\
Laboratoire d'Hydrologie Mathématique

\uline{Résumé de la thèse} :

L'étang-solaire est un capteur solaire constitué d'un bassin peu profond, partiellement transparent au rayonnement solaire, et dont les pertes thermiques sont réduites en éliminant les phénomènes convectifs par un gra-\\
dient salin et stabilisateur.

Plusieurs de ces bassins ont déjà été étudiés dans le monde. Un bassin expérimental a été construit à Palavas et exploité en collaboration, par l'Université des Sciences et Techniques du Languedoc, le Centre National pour l'Exploitation des Océans et la Société ELF.

Deux campagnes de mesures, de mars 1977 à décembre 1977 et de mars 1978 à mai 1979, ont permis de résoudre les problèmes technologiques de mise en place et maintien de la stratification saline, et de suivre l'évolution thermique du milieu.

Cette expérimentation a rendu possible la mise au point d'un modèle d'évolution d'un tel bassin à partir de la connaissance des données géographiques, météorologiques et techniques du bassin.

Des tests de validité ont montré l'exactitude du modèle. Des simulations ont permis la définition des rendements du système sous diverses conditions du fonctionnement (\textit{température d'extraction,quantité de chaleur extraite}).

\uline{Durée approximative des recherches ayant abouti à la thèse} : 2 ans et demi.

\uline{Composition du jury proposé} :
\begin{tabular}{lll}
MM.:& BEN AIM & R.\\
& BOCQUILLON & Cl.\\
& HYACINTHE & J-L.\\
& DE CHABERT & F.\\
& ROUSSEL & A.\\
& PERRUT & M.\\
& RONC & M.
\end{tabular}

Le Secrétaire Général,

\unsure{E S}

\uline{E. S I A U}
\mspage{60}
